\section{Methods}
\label{sec:methods}

Every method takes a pool and returns a ranked list of \emph{axes}. An axis is a
set of failing routes the method claims belong together. Scoring compares the
method's first ten distinct axes with each gap's failing routes
(\S\ref{sec:scoring}). Code paths below are relative to
\code{docs/research/experiments/f2d-dev-simlingo/} unless they start with
\code{src/}.

\subsection{The sprint framework}
\label{sec:sprint}

The search ran as a \emph{discovery sprint} (\code{docs/prd-discovery-sprint.md}).
Its rules are these.
\begin{enumerate}
  \item A method is proposed and its row is written in the variant log
        (\code{VARIANTS.md}) before its score is read. Chris approved each batch
        before it ran.
  \item One CPU pass (\code{cache\_pool.py}) turns a pool's shards into a single
        per-route summary file, so a method runs in minutes and is scored by one
        command (\code{sprint\_score.py}).
  \item Every scored configuration counts as a \emph{try}, failures and
        cancelled runs included. A rescore of a logged configuration under a new
        rule is not a try. Every dev table states the running try count, and
        every dev $p$ value is exploratory.
  \item At most three finalists are frozen and tagged, then scored once on each
        held-out model. One is declared primary in the tag, so the headline
        needs no correction for multiple comparisons.
\end{enumerate}
% measured: try counts 134 (A2), 188 (A3), 202 (A4), 216 (unpooled), 219 (views tag); docs/prereg-fail2drive.md §8-§10, docs/prereg-views.md, VARIANTS.md
The try count was 134 when amendment A2 was written, 188 at the freeze, 202 when
the search stopped, and is 219 once the views pre-registration's new rows are
counted. Two people's agents searched in parallel: Chris's agents wrote the
\emph{m} family and Jerry's explore loop wrote the \emph{j} family. Both logged
into the same file and scored with the same scorer.

\subsection{Shared inputs}
\label{sec:inputs}

\paragraph{Failure step.}
For a failing route, the failure step is the step of its earliest timed
infraction (\code{src/harness/failure\_time.py}). Collisions carry a logged step.
Other infractions carry a location, and are timed at the first step within
0.5\,m of the ego's nearest approach to it.
% design: 0.5 m tolerance (src/harness/failure_time.py, LOCATED_TOLERANCE_M)
A failing route with no timed infraction, and every passing route, uses its end.

\paragraph{Summaries.}
For every route, site and layer, \code{cache\_pool.py} stores three means of the
per-frame activations: the \emph{whole route}, the \emph{approach window} (the
last $W$ captured frames before the failure step), and the \emph{early route}
(the first $W$ frames).
% design: W = 6 by default; windows of 3, 10, 20, 30, 50 and 100 frames tried (VARIANTS.md batches 2-4)
The default is $W=6$ frames. Since frames are 0.25\,s apart, 6 frames is about
1.5\,s and 50 frames about 12.5\,s.
% measured: 0.25 s spacing (WHAT_BROKE.md 2026-11-16)
Rows are written in salted-hash order, never in route-id order
(\S\ref{sec:whatbroke}, H25). A per-frame file (\code{frames\_pool.py}) keeps the
last 50 frames of a few layers for the methods that model frame sequences.

\paragraph{Layers.}
``L8'' means layer 8 of SimLingo's \code{driving\_queries} site. When several
layers are concatenated, each is standardised and divided by
$\sqrt{n_\text{layers}}$ so the layers weigh equally.

\subsection{The core recipe, m03}
\label{sec:m03}

Most methods in both families are variations on one recipe, m03
(\code{cache\_methods.py}).
\begin{enumerate}
  \item Take one summary at one site and layer for every route, passing and
        failing: a matrix with one row per route.
  \item Standardise each column over all routes.
  \item Fit PCA on all routes and keep 32 components. Standardise the scores.
        % design: 32 components (cache_methods.py --dims default; VARIANTS.md batch 1)
  \item Keep only the failing rows. Build a Ward linkage tree on them.
  \item Cut the tree into $k=2$, then $3$, then $5$ clusters. Within each cut,
        list clusters largest first. Under A2, a cluster already listed at an
        earlier cut is not listed again; if fewer than 10 distinct clusters
        result, keep cutting at $k=6, 7, \dots$ until there are 10
        (\code{cache\_methods.axes\_from}).
  \item The listed clusters, in order, are the method's axes.
\end{enumerate}
Only the pass/fail outcome is read. The fit runs on the pool being scored, so no
model is trained on Fail2Drive in the sense its authors forbid.

\subsection{The m family}
\label{sec:mfamily}

Table~\ref{tab:mfamily} lists the m methods. Each row changes one thing about
m03 or combines earlier rows. The settings are those logged in
\code{VARIANTS.md}; each id there links to its job and script.

\paragraph{Pooling.}
Several m methods merge the ranked axes of two or three other methods
(\code{pool\_axes.py}). The merge is round-robin in a fixed order: each source's
first axis, then each source's second, and so on. An axis whose Jaccard with one
already kept is at least 0.8 is skipped. It stops at 10 axes. The order is fixed
before scoring and nothing reads the key.
% design: 0.8 skip threshold, stop at 10 (pool_axes.py)

\paragraph{The two finalists.}
Both are frozen at tag \code{f2d-freeze-1}.
\begin{itemize}
  \item \textbf{m23 (primary).} Source 1 is m03 on the 50-frame approach window
        at L8 (this is m15). Source 2 is m03 on the whole-route summary at L8.
        The two ranked lists are pooled in that order. The idea is that the
        approach window carries failures that happen in the last seconds, and
        the whole-route summary carries failures that unfold over the route.
  \item \textbf{m40 (secondary).} Source 1 is m03 on the 50-frame window at L8
        with 64 PCA components. Source 2 is m03 on the whole-route summary at
        L20. Pooled in that order.
\end{itemize}

\begin{table}[t]
\centering
\small
\begin{tabular}{@{}lp{0.78\linewidth}@{}}
\toprule
Id & What changes from m03 \\
\midrule
m01, m02 & Error-aware B7: B7's sparse-dictionary latents ranked by their
           enrichment in failing routes; m02 on window frames only \\
m04      & A balanced logistic pass/fail probe on the PCA scores picks the
           components that separate failure; cluster in that subspace \\
m06      & m03 plus B3-behaviour's kinematic features, equal block weight \\
m07      & Window mean minus the route's own early mean \\
m05      & Sweep of m03 over 24 layers of two sites plus \code{vision\_bridge} \\
m08, m16 & Several layers concatenated (L5, L8, L20) \\
m11, m15, m24 & Longer approach windows: 3, 10, 20, 30, 50, 100 frames \\
m12, m17 & Average, complete or spectral linkage in place of Ward \\
m13      & Other cut schedules (4 and 6; 10) \\
m19      & Layers chosen by pass/fail probe AUROC (uninformative; see
           \S\ref{sec:whatbroke}) \\
m09, m10, m20, m23, m28 & Pools of other methods' axes (\code{pool\_axes.py}) \\
m22, m27 & Two summaries or two sites concatenated into one embedding \\
m25      & Whole-route summary at other layers \\
m26      & Window minus whole-route mean \\
m29--m32 & Per-frame summaries: several time scales, PCA trend, stopped versus
           moving frames, mean and spread \\
m33--m36 & j09 at other layers, beside the window mean, or on several layers \\
m38      & j05's background subtraction on the L8 50-frame window \\
m39, m40 & PCA width 16 or 64 \\
\bottomrule
\end{tabular}
\caption{The m family. Round-2 tries (on SimLingo and TransFuser++ pooled) are
logged in \code{VARIANTS.md} under ``Round 2'' and follow the same pattern.}
\label{tab:mfamily}
\end{table}

\subsection{The j family}
\label{sec:jfamily}

The j methods came from Jerry's explore loop. Each asks whether a different
treatment of the activation geometry finds groups that m03 misses. Table
\ref{tab:jfamily} lists them. All read only the pass/fail outcome, and all end
in \code{axes\_from} or the same largest-first ordering unless stated.

\begin{table}[t]
\centering
\small
\begin{tabular}{@{}lp{0.82\linewidth}@{}}
\toprule
Id & Method \\
\midrule
j01 & m03 without standardising: PCA scores keep their eigenvalue scale; each
      layer is divided by its root-mean-square instead \\
j02 & Contrastive PCA: directions where failing routes vary more than passing
      ones, after a 32-component unwhitened PCA \\
j03 & m03 embedding cut finer (8, 12, 16 clusters), keeping only clusters in a
      size band, ranked by cohesion \\
j04 & Whiten with the passing routes' covariance, keep the coordinates with the
      largest failure second moment \\
j05 & Unwhitened PCA; subtract from each failing route the mean of its 10
      nearest passing routes; cluster the residuals \\
j06 & j01 after removing a least-squares prediction of the activations from the
      behaviour features \\
j07 & Split failing routes into two behaviour strata, then j01 within each \\
j08 & Candidate clusters from several views (behaviour and three layers), ranked
      by their stability under subsampling of the failing routes \\
j09 & A diagonal Gaussian hidden Markov model (a few hidden states, each with its
      own Gaussian, plus transition probabilities) fit on every route's last 50
      frames at L8 after PCA to 16; failing routes clustered on state occupancy
      and transitions \\
j10 & j09's model, with each route embedded by its Fisher score on the emission
      means \\
j11 & j05's residuals scaled to unit length \\
j12 & Residual from a Gaussian mixture fit on passing routes only \\
j13 & j05's residual minus its projection on the local principal directions of
      its 10 nearest passing routes \\
j14 & Window score minus the mean window score of the nearest passing routes in
      early-route space \\
j15 & j05's residual clustered by HDBSCAN~\cite{campello2013hdbscan,mcinnes2017hdbscan},
      clusters ranked by stability \\
\bottomrule
\end{tabular}
\caption{The j family. Scripts are \code{j01\_unwhite.py} to \code{j15\_hdbscan.py};
settings and commits per run are in \code{VARIANTS.md}.}
\label{tab:jfamily}
\end{table}

\subsection{Rivals}
\label{sec:rivals}

Each rival ranks its own axes and is scored exactly like a method.
\begin{description}
  \item[B1, chance.] Random axes of the same sizes as the method's own, drawn
        from failing routes only, averaged over 50 draws.
        % design: 50 draws (sprint_score.py, chance)
  \item[B3-behaviour, the primary rival.] Built only from signals a fleet logs
        (\code{b3\_behaviour.py}, tag \code{f2d-sprint-a1}). Per route: in the
        approach window (the last 6 captured frames, about 1.5\,s) the mean, minimum,
        maximum and final speed, the mean acceleration and the largest
        deceleration; over the whole route the mean speed, the fraction of frames
        below 0.5\,m/s and the log of the step count; and one count per infraction
        type seen in the pool. Features are standardised over all routes; failing
        routes are Ward-clustered and cut as in m03. Excluded: distances to
        objects, any scenario or route field, and the action error against the
        rule-based expert, since a fleet runs no expert beside the car.
        Steering and braking are not in the shards.
  \item[B3-infraction, a reported rival.] One axis per infraction type: the
        failing routes with at least one of it, plus one axis for failing routes
        with no recorded infraction, largest first (\code{b3\_infractions.py}).
        It is what a fleet team does first. It is compared with every method but
        has no pass bar and is not part of the best-rival comparison (A2).
  \item[B7-pool.] Activation statistics with no outcome labels, after the
        feature statistics of Dr.~VLA~\cite{swann2026sparse}: per-latent episode
        coverage, onset count, peak activation and run length of the latents of a
        TopK sparse autoencoder (\code{train\_saes.py}; settings in
        \code{PRD.md} \S12.2), used as a gap signal. That use is our
        construction, not Dr.~VLA's (\code{docs/steps/b7-gap-signal.md}). One
        dictionary only, the primary \code{L10\_k64\_s0}, so the rival is not the
        best of a grid.
  \item[B4.] The enumerated factor list. It decides the B4-named split
        (\S\ref{sec:fail2drive}) and is reported as a rival.
  \item[B2 with labels.] An oracle ceiling that reads scenario fields; never a
        rival.
\end{description}
Full B3, which also read the simulator's exact distances to nearby objects, was
retired at A1 because no deployed car has them.

\subsection{Scoring under A1 and A2}
\label{sec:scoring}

The scorer is \code{sprint\_score.py}; the harness definitions are in
\code{src/harness/scoring.py}.

\paragraph{Own top ten.}
No axis is chosen with the key. Each method's axes are taken in its own order.
Exact repeats are dropped (A2) and the first $M=10$ distinct axes count.
% design: M = 10 (docs/prereg-fail2drive.md §3, §7; src/harness/scoring.py BUDGET_M)

\paragraph{Primary: mean best Jaccard.}
For a gap $g$ with failing routes $G$ and axes $A_1,\dots,A_{10}$,
\begin{equation}
  \text{best}(g) = \max_{i \le 10} J(G, A_i), \qquad
  J(G, A) = \frac{|G \cap A|}{|G \cup A|}.
\end{equation}
The score is the mean of $\text{best}(g)$ over the gaps, each weighted equally.
Two edge cases are fixed on purpose. First, $J(\emptyset, \emptyset) = 0$, not 1;
otherwise a method could score perfectly on an empty gap by emitting nothing.
Second, a method with no axes scores 0.

\paragraph{Beside it.}
Recall at Jaccard 0.5: one Hungarian assignment~\cite{kuhn1955hungarian} of gaps
to axes on the Jaccard matrix, counting matched pairs with $J \ge 0.5$, divided
by the number of gaps.
% design: threshold 0.5 (src/harness/scoring.py MATCH_THRESHOLD)
Decoy hit rate: the share of axes at least half of whose routes are slice
routes of one decoy type. Precision, where the harness reports it, is the number
of matched axes divided by the axis budget $M$, never by the number of axes the
method emitted; otherwise a method that emits one safe axis would beat a full
directive.

\paragraph{Rank-averaged best match.}
The views pre-registration makes this the primary metric
(\code{docs/research/experiments/f2d-unpooled/view\_score.py}). For a gap it is
the mean over $k=1,\dots,10$ of the best Jaccard among the method's top $k$
axes. It rewards a method that ranks the matching axis early. Top-10 best match
is reported beside it.

\paragraph{Pair bootstrap.}
Routes from the same Fail2Drive pair share a road, and routes from different
seeds of one pair share both twins, so routes are not independent. The resampling
unit is the pair, carrying all its routes on every seed (A1).
\begin{enumerate}
  \item List the distinct pairs in the pool. The pair id must not carry the seed
        (\S\ref{sec:h21b}).
  \item In each of 2{,}000 replicates, draw that many pairs with replacement.
        % design: 2,000 replicates (docs/prereg-fail2drive.md §7; docs/prereg-views.md)
        Each route gets weight equal to the number of times its pair was drawn.
  \item Recompute every score with weighted Jaccard:
        $J_w(G,A) = \sum_{r \in G \cap A} w_r \big/ \sum_{r \in G \cup A} w_r$,
        with 0 when the union weighs 0.
  \item A contrast is method minus rival in the same replicate. Report the point
        difference, the 2.5th and 97.5th percentiles, and the one-sided $p$, the
        share of replicates where the difference is at most 0.
\end{enumerate}
When a contrast is against ``the best rival'' or ``the best new view'', the best
is re-chosen inside every replicate, so the interval pays for the choice. Pooled
results across models are a fixed mean over units with no random effect.

\subsection{Running a frozen method on another model}
\label{sec:transfer}

The finalists are written in SimLingo's terms, so \code{transfer-spec.md}
(binding from \code{f2d-freeze-1}) maps them.
\begin{itemize}
  \item \textbf{Site:} \code{driving\_queries} on every model, the representation
        each controller acts on.
  \item \textbf{Layer by relative depth:} SimLingo layer $L$ is fraction
        $f = L/23$; on a site of depth $D$ the layer is
        $\lfloor f(D-1) + 0.5 \rfloor$. On TransFuser++ ($D=6$), m23 runs at
        layer 2 for both views; m40 at layer 2 (window) and 4 (whole route).
  \item \textbf{Window in frames:} unchanged, since all three models capture
        every 5 steps.
  \item \textbf{Everything else unchanged}, fit on the held-out model's own pool.
  \item \textbf{B7-pool's sparsity is kept, not its $k$:} $k = \text{round}(64
        \times \text{width} / 896)$, so $k=18$ at width 256.
        % design: k rule fixed at A3 (transfer-spec.md item 5)
  \item \textbf{UniAD fallback:} if UniAD's shards lack ego speed or step,
        B3-behaviour uses the infraction counts and the route's step count, and
        this is stated beside every UniAD result (A3).
\end{itemize}
